\documentclass[12pt]{article}
\usepackage[utf8]{inputenc}
\usepackage[T1]{fontenc}
\usepackage{hyperref}
\title{Bannoni Language Resources for K-12 Curriculum Development}
\author{Compiled by Claude Sonnet 5.5 (claude-sonnet-5-5) for Andr\'as Kornai}
\date{}
\begin{document}
\maketitle

\section{Introduction}
Bannoni (ISO 639-3 bcm; Glottocode bann1247), also spelled Banoni and called Tsunari by its speakers, is an Austronesian language of Empress Augusta Bay in southwest Bougainville, Papua New Guinea. Speaker estimates in the sources found conflict (about 1,000 to 4,500). It is closely related to Piva. Dialects include Amun, Torokina and Mabes. Documentation is mostly archival and linguistic.

\section{Primary Educational Resources}
No primers, textbooks or school programmes were found in search results.

\section{Linguistic and Reference Materials}
\begin{itemize}
\item Wikipedia, Bannoni language. \url{https://en.wikipedia.org/wiki/Bannoni_language}
\item PARADISEC field notes series NSBBAN (Capell papers, Banoni and dialects, including a vocabulary list and notes on orthography). \url{https://paradisec.org.au/fieldnotes/NSBBAN.htm}
\item PARADISEC catalog entry. \url{https://catalog.paradisec.org.au/repository/BP3/006}
\item Surrey Morphology Group, Northwest Solomonic bibliography (Piva/Banoni). \url{https://www.smg.surrey.ac.uk/languages/northwest-solomonic/bibliography/pivabanoni/}
\end{itemize}

\section{Religious and Community Resources}
\begin{itemize}
\item Joshua Project, Banoni in Papua New Guinea. \url{https://joshuaproject.net/people_groups/10659}
\end{itemize}

\section{Digital and Audio Resources}
The PARADISEC archive holds recordings and papers; consult the catalog entry above for access conditions.

\section{Curriculum Development Context}
The language has a written form according to Joshua Project, but no school use was found. A curriculum would start from the Capell orthography notes and archival materials, with community and Bougainville education authorities.

\end{document}